\section{A transfer criterion requiring only positive error in mean}
\label{sec:mean-transfer}

This section gives a self-contained strengthening of the analytic interface used in
the repository.  It does not establish a new zero-free region for a new
$L$-function: the arithmetic estimate in the application below remains to be
proved.  The contribution relative to the incoming reports is to replace the
pointwise arithmetic target by a positive-part estimate in mean, and to state
exactly how nonuniform rank and denominator losses can feed that estimate.
The analytic argument is an adaptation of classical Landau theory, not a claim
that the underlying analytic principle is new to the literature.

For a real number $u$, write $u_+=\max(u,0)$ and $u_-=\max(-u,0)$.

\begin{theorem}[Positive-part Mellin criterion]\label{thm:positive-mellin}
Let $\beta>0$, and let $R:[1,\infty)\to\mathbb R$ be measurable and locally
integrable, with $|R(x)|\le Cx^A$ for all sufficiently large $x$, for some
finite constants $A,C$. Suppose
\[
 F(s)=\int_1^\infty R(x)x^{-s}\,dx
\]
initially defined for $\Re s>A+1$, has a meromorphic continuation to
$\Re s>\beta$ with at most simple poles.  At every real $u>\beta$, suppose
the residue of $F$ is a nonnegative real number, with residue zero at a
regular point.  If
\begin{equation}\label{eq:positive-integrability}
 \int_1^\infty R(x)_+x^{-\sigma}\,dx<\infty
 \qquad\hbox{for every real }\sigma>\beta,
\end{equation}
then $F$ is holomorphic throughout $\Re s>\beta$. Moreover
\[
 \int_1^\infty |R(x)|x^{-\sigma}\,dx<\infty
 \qquad(\sigma>\beta).
\]
\end{theorem}

\begin{proof}
Put $H(s)=\int_1^\infty R(x)_+x^{-s}\,dx$. Condition
\eqref{eq:positive-integrability}, with a slightly smaller real exponent to
dominate derivatives, makes $H$ holomorphic on $\Re s>\beta$. Far to the
right,
\[
 G(s):=H(s)-F(s)=\int_1^\infty R(x)_-x^{-s}\,dx.
\]
The right-hand side is the Mellin transform of a nonnegative function. Let
$a\in[-\infty,\infty)$ be its abscissa of convergence; the polynomial growth
hypothesis makes $a$ bounded above. It is enough to prove $a\le\beta$.

Suppose $a>\beta$. The meromorphic function $H-F$ continues $G$ to a
neighborhood of $a$. It agrees with $G$ in $\Re s>a$ by the identity theorem.
Its residue at $a$ equals $-\operatorname{Res}_{s=a}F(s)\le0$. On the other
hand, $G(\sigma)\ge0$ for real $\sigma>a$, so
\[
 \operatorname{Res}_{s=a}(H-F)
 =\lim_{\sigma\downarrow a}(\sigma-a)G(\sigma)\ge0.
\]
Thus the residue is zero. Because the possible pole is at most simple,
$H-F$ is holomorphic at $a$. This contradicts Landau's theorem: a Mellin
transform of a nonnegative locally integrable function, equivalently its
Laplace transform after $x=e^v$, cannot be holomorphic at its finite abscissa
of convergence. Hence $a\le\beta$.

Both positive and negative parts now have absolutely convergent transforms
for every real $\sigma>\beta$. Their difference is holomorphic there and,
by analytic continuation, equals $F$. Their sum proves the final assertion.
\end{proof}

The exact source interface for the Landau step is
\path{Lib/Analytic/Landau.lean}, especially
\path{exists_Ioi_subset_integrableExpSet_of_analyticAt}
(lines 160--176) and
\path{exists_Ioi_subset_integrableExpSet}
(lines 182--210) at the pinned commit.  The repository's
\texttt{IsPolar.landau} (lines 217--259) instead subtracts a fixed
exponential majorant. The present proof subtracts the holomorphic transform
of the whole positive part \cite{landau}.

\begin{corollary}[Dyadic and square-mean versions]\label{cor:mean-versions}
Under the analytic hypotheses of Theorem~\ref{thm:positive-mellin}, either
of the following conditions suffices to exclude all poles with $\Re s>\beta$:
\begin{align}
 \int_X^{2X}R(x)_+\,dx&\ll_\varepsilon X^{\beta+\varepsilon},
       &&\hbox{for every }\varepsilon>0,\label{eq:mean-lone}\\
 \int_X^{2X}R(x)_+^2\,dx&\ll_\varepsilon X^{2\beta-1+\varepsilon},
       &&\hbox{for every }\varepsilon>0.\label{eq:mean-ltwo}
\end{align}
Here it is enough that the estimates hold for all sufficiently large
dyadic $X=2^k$. In fact \eqref{eq:mean-lone} is equivalent to
\eqref{eq:positive-integrability}.
\end{corollary}

\begin{proof}
For $\sigma>\beta$, choose $0<\varepsilon<\sigma-\beta$ and sum
\[
 \int_{2^k}^{2^{k+1}} R(x)_+x^{-\sigma}\,dx
 \ll_\varepsilon 2^{k(\beta+\varepsilon-\sigma)}.
\]
The geometric series converges. Conversely, integrability at
$\sigma=\beta+\varepsilon$ bounds a dyadic block by a constant times
$X^{\beta+\varepsilon}$. Finally Cauchy--Schwarz gives
\[
 \int_X^{2X}R(x)_+\,dx
 \le X^{1/2}\left(\int_X^{2X}R(x)_+^2\,dx\right)^{1/2},
\]
which takes \eqref{eq:mean-ltwo} to \eqref{eq:mean-lone}.
\end{proof}

The integral condition really allows positive spikes that no eventual
pointwise bound allows. For example, let $R$ be the indicator of the disjoint
intervals $[2^k,2^k+2^{-2k}]$, $k\ge1$. Its Mellin transform is holomorphic
on $\Re s>-2$, and its positive part satisfies the criterion for
$\beta=1/2$, but $R(x)\le x^{-1/2+\varepsilon}$ fails arbitrarily far out
for every $0<\varepsilon<1/2$.

\subsection{An exact averaged determinant certificate}

The source proof compares a lower bound on the same covolume plus twice an
index with an upper bound obtained from many short vectors. Its scalar
content can be retained while allowing the fraction of available vectors to
vary strongly with $x$.
The scalar rearrangement already appears in the predecessors; its use
below with an averaged inverse-rank loss is the additional interface.

\begin{proposition}[Rank-sensitive loss]\label{prop:rank-loss}
At each $x$ suppose $d>0$, $0<n\le d$, $M>0$, $H\ge0$, and $L,U,I\ge0$.
Suppose a real number $V$ satisfies
\[
 d(H-L)\le V+2dI,
 \qquad
 V\le n\bigl(H-MR(x)+U\bigr).
\]
Writing $\alpha=n/d$, one has the exact scalar implication
\begin{align}
 R(x)_+
 &\le\frac{\bigl(L+\alpha U+2I-(1-\alpha)H\bigr)_+}{\alpha M}
 \label{eq:rank-exact}\\
 &\le\frac{1}{M}\left(U+\frac{L+2I}{\alpha}\right).
 \label{eq:rank-coarse}
\end{align}
\end{proposition}

\begin{proof}
Combine the two inequalities, divide by $d$, and rearrange to obtain
\[
 \alpha M R(x)\le L+\alpha U+2I-(1-\alpha)H.
\]
Taking positive parts and dividing by $\alpha M>0$ proves
\eqref{eq:rank-exact}. Since $H\ge0$, $\alpha\le1$, and
$L+\alpha U+2I\ge0$, dropping the nonpositive term gives
\eqref{eq:rank-coarse}.
\end{proof}

Thus one must first combine the bounds for each $x$, and then average the
combined loss. Let
\[
 A(x)=U(x)+\frac{L(x)+2I(x)}{\alpha(x)}.
\]
If $M(x)\ge c x^\theta$ with $0<\theta<1$, and for every
$\varepsilon>0$
\begin{equation}\label{eq:mean-arithmetic-target}
 \int_X^{2X} A(x)\,dx\ll_\varepsilon X^{1+\varepsilon},
\end{equation}
then \eqref{eq:rank-coarse} implies
\[
 \int_X^{2X}R(x)_+\,dx
 \ll_\varepsilon X^{1-\theta+\varepsilon}.
\]
If $R$ has the analytic properties of Theorem~\ref{thm:positive-mellin},
this excludes poles with $\Re s>1-\theta$. The factor $M$ may itself
be a product of the number of lifts and the normalization of the prime
statistic; keeping that normalization is essential when applying the formula.

There is no requirement of a uniform positive lower bound for $\alpha$.
For example, conjugate exponents $p,q>1$ with $1/p+1/q=1$ give the
sufficient conditions
\begin{align*}
 \int_X^{2X}U(x)\,dx&\ll_\varepsilon X^{1+\varepsilon},\\
 \int_X^{2X}(L(x)+2I(x))^p\,dx&\ll_\varepsilon X^{1+\varepsilon},\\
 \int_X^{2X}\alpha(x)^{-q}\,dx&\ll_\varepsilon X^{1+\varepsilon}.
\end{align*}
H\"older's inequality proves \eqref{eq:mean-arithmetic-target}. The
case $p=q=2$ gives a particularly concrete square-mean test. Bounding
the averages of $L,I,\alpha$ separately is insufficient: on a probability
space, let $L=1$, and let $\alpha=\delta^2$ with probability $\delta$
and $\alpha=1$ otherwise. Then $\mathbb E L=1$ and
$\mathbb E\alpha=1-\delta+\delta^3\to1$, but
$\mathbb E(L/\alpha)=1-\delta+\delta^{-1}\to\infty$.

\subsection{Application to the nonabelian cubic target}

Let $K=\mathbb Q(\sqrt{-23})$, let $H$ be its Hilbert class field, and
let $\mathcal R_H$ be the earlier report's principal-ideal remainder
\[
 \mathcal R_H(x)=
 3\sum_{\substack{N\mathfrak a\le x\\{}[\mathfrak a]=1}}
 \Lambda_K(\mathfrak a)\left(\frac1{N\mathfrak a}-\frac1x\right)
 -\log x+1+\gamma_H.
\]
Its Mellin transform is
\[
 \frac{-\zeta_H'(s)/\zeta_H(s)}{s(s-1)}
 -\frac1{(s-1)^2}+\frac{1+\gamma_H}{s-1}.
\]
The singularity at $s=1$ is removable. A nontrivial zero $\rho$ of
multiplicity $m_\rho$ gives residue
$m_\rho/[\rho(1-\rho)]$, which is nonnegative at real
$0<\rho<1$. The coefficient bound
$0\le\Lambda_H(n)\le6\Lambda(n)$ gives the required elementary growth.
Consequently either one of
\begin{align*}
 \int_X^{2X}\mathcal R_H(x)_+\,dx
    &\ll_\varepsilon X^{1/2+\varepsilon},\\
 \int_X^{2X}\mathcal R_H(x)_+^2\,dx
    &\ll_\varepsilon X^\varepsilon
\end{align*}
for every $\varepsilon>0$ would prove RH for $\zeta_H$ and hence for the
degree-two class-character $L$-function and the cubic Dedekind zeta function.
These require a weaker form of arithmetic control than the earlier eventual
pointwise bound; for this fixed remainder, successful control would still
be equivalent to the same zero-location assertion. They remain unproved here. The new calculation identifies what an
arithmetic construction in mean would have to deliver; it does not supply
that construction.

\subsection{Signed moving averages also suffice, if every starting point is controlled}

There is a different useful relaxation in which cancellation is allowed but
the starting point of the averaging interval must vary continuously. This
must be distinguished from checking only the fixed dyadic partition.

\begin{theorem}[Signed moving-window criterion]\label{thm:moving-window}
Assume the analytic hypotheses of Theorem~\ref{thm:positive-mellin}, and
fix $h>1$. If, for some $C\ge0$,
\begin{equation}\label{eq:moving-signed}
 \int_x^{hx}R(t)\,dt\le Cx^\beta
 \qquad\hbox{for every sufficiently large real }x,
\end{equation}
then $F$ is holomorphic in $\Re s>\beta$.
It also suffices, for a desired boundary $\beta$, to prove
\eqref{eq:moving-signed} with exponent $\beta+\varepsilon$ for every
$\varepsilon>0$, with constants allowed to depend on $\varepsilon$.
\end{theorem}

\begin{proof}
Set $T_hR(x)=x^{-1}\int_x^{hx}R(t)\,dt$. Fubini's theorem, initially
far to the right, gives
\begin{align*}
 F_{T_hR}(s)
 &=\frac{h^s-1}{s}F(s)
 -\frac1s\int_1^h R(t)\bigl(h^s t^{-s}-1\bigr)\,dt.
\end{align*}
Indeed the inner integral in the reversed order is
$\int_{\max(1,t/h)}^t x^{-s-1}\,dx$; splitting the outer integral
at $t=h$ proves the formula. The finite-interval correction is entire:
the apparent singularity at $s=0$ is removable. The multiplier
$A_h(s)=(h^s-1)/s$ is entire, is nonzero on $\Re s>0$, and is
strictly positive for real $s>0$. Thus the poles of $F_{T_hR}$ in
$\Re s>\beta$ have exactly the poles of $F$, and their real residues
are nonnegative.

The estimate \eqref{eq:moving-signed} gives
$T_hR(x)\le Cx^{\beta-1}$ eventually. Since $R$ is locally integrable,
$T_hR$ is continuous and bounded on each compact interval. Enlarge $C$ to
$K$ so that $Kx^{\beta-1}-T_hR(x)\ge0$ for all $x\ge1$.
Its transform is
$K/(s-\beta)-F_{T_hR}(s)$. At real points greater than $\beta$
its residues are nonpositive. The same Landau abscissa argument as above
therefore shows that $F_{T_hR}$ has no pole to the right of $\beta$.
Since $A_h$ has no zero there, neither does $F$ have a pole. Apply the
fixed-exponent result for each $\beta+\varepsilon$ to get the last claim.
\end{proof}

In particular, for the cubic-field target one may instead aim for
$\int_x^{2x}\mathcal R_H(t)\,dt\le C_\varepsilon
x^{1/2+\varepsilon}$ for \emph{every} sufficiently large real $x$.
This allows signed cancellation.  No corresponding arithmetic estimate is
proved here. The positive-part criterion is more robust when the available
information is restricted to the sampled dyadic intervals.

\subsection{Why signed averages on the dyadic grid are insufficient}

The positive part in the criterion cannot simply be removed. Fix
$0<\beta<a<1$, put $\tau=2\pi/\log2$, and set
\[
 R(x)=\frac{d}{dx}\bigl(x^a\sin(\tau\log x)\bigr)
 =x^{a-1}\bigl(a\sin(\tau\log x)+\tau\cos(\tau\log x)\bigr).
\]
For every integer $k\ge0$,
\[
 \int_{2^k}^{2^{k+1}}R(x)\,dx=0.
\]
Nevertheless
\[
 F(s)=\int_1^\infty R(x)x^{-s}\,dx
 =\frac{\tau s}{(s-a)^2+\tau^2}
 \qquad(\Re s>a),
\]
which has nonzero simple poles at $a\pm i\tau$, to the right of
$\beta$, and no real poles. All other analytic hypotheses of
Theorem~\ref{thm:positive-mellin} hold. Moreover
\[
 \int_{2^k}^{2^{k+1}}R(x)_+\,dx=C_a\,2^{ka}
 \quad\hbox{with }C_a>0.
\]
This follows by substituting $x=2^k e^u$, $0\le u\le\log2$.
Thus cancellation in signed bin averages can hide precisely the poles
that the arithmetic transfer is intended to exclude. This example concerns
signed averages on the fixed dyadic partition; it does not assert failure of
every translated smoothing condition.
